% ======================================================================
% SECTION 6: AUTONOMOUS REGULATORY AFFAIRS AND SUBMISSIONS
% File: sections/regulatory_affairs.tex
% ======================================================================

[PROCESSING INSTRUCTION: Generate this section (approximately 3-4 pages) covering autonomous
regulatory affairs operations. Use the following source files:

SOURCE FILES TO PROCESS:
- sponsor/input\_files/sponsor\_02\_preclinical\_regulatory.md (IND/CTA mechanics, accelerated
  pathways, regulatory interactions, global filing strategies)
- sponsor/input\_files/sponsor\_06\_cmc\_submission.md (eCTD v4.0, CDISC, submission readiness,
  accelerated mechanisms, global markets)
- sponsor/input\_files/org\_01\_executive\_regulatory.md (sponsor regulatory accountability)
- sponsor/input\_files/org\_07\_cro\_split\_references.md (sponsor vs. CRO regulatory split)
- national-platform/new\_paper/final\_paper/sections/regulatory\_landscape.tex (regulatory landscape)
- national-platform/new\_paper/final\_paper/sections/cfr312\_adaptation.tex (adapted 21 CFR 312)
- national-platform/new\_paper/final\_paper/sections/ich\_e6r3\_adaptation.tex (adapted ICH E6(R3))

SUBSECTION STRUCTURE:

\subsection{Regulatory Agent (regulatory\_agent)}
- IND lifecycle management: initial filing, amendments, annual reports
- Authority correspondence package generation
- Multi-jurisdiction coordination (FDA, EMA, PMDA, Health Canada, TGA)
- Maps to: Regulatory Affairs team

\subsection{IND and CTA Filing Automation}
- U.S. IND: 30-day review window, serial submissions
- EU CTIS transition (deadline 31 January 2025 past, ongoing management)
- UK CTIMP reform (April 2026)
- eCTD v4.0 (September 2024) document assembly

\subsection{Accelerated Pathway Navigation}
- FDA accelerated mechanisms: Breakthrough, Fast Track, Priority Review, Accelerated Approval
- Project Orbis concurrent multi-country review
- RTOR real-time oncology review
- EMA: PRIME, accelerated assessment, conditional MA
- PMDA: SAKIGAKE, conditional early approval
- Global markets: Health Canada NOC/c, Swissmedic fast-track, TGA provisional, NMPA conditional

\subsection{Regulatory Intelligence and Horizon Scanning}
- Automated monitoring of regulatory guidance updates
- Impact assessment on active programs
- Deadline tracking and compliance calendars
- Integration with regulatory\_agent decision logic

PYTHON SCRIPT REQUIREMENTS (to be generated by Claude Code Opus 4.6 in the future):
- File: sponsor/template/scripts/regulatory\_agent.py
  Purpose: IND lifecycle manager. Track IND status, generate amendments, produce
  annual reports (21 CFR 312.33), manage authority correspondence. Include
  multi-jurisdiction support for FDA/EMA/PMDA/HC/TGA filings.
- File: sponsor/template/scripts/ind\_filing\_engine.py
  Purpose: Automated IND/CTA package assembly. Parse protocol, IB, CMC data,
  nonclinical summaries into eCTD v4.0 structure. Generate Module 1-5 documents.
  Include CTIS-specific requirements for EU submissions.
- File: sponsor/template/scripts/accelerated\_pathway\_advisor.py
  Purpose: Evaluate program eligibility for accelerated pathways. Score against
  criteria for Breakthrough, Fast Track, PRIME, Orbis, RTOR. Generate
  designation request packages. Input: indication, unmet need data, preliminary
  efficacy. Output: pathway recommendation with supporting rationale.
- File: sponsor/template/scripts/regulatory\_intelligence\_scanner.py
  Purpose: Monitor regulatory guidance repositories (FDA, EMA, ICH) for updates.
  Parse guidance documents, assess impact on active programs, generate alerts.
  Include deadline tracking for E2B(R3) April 2026 and other time-sensitive items.
- All scripts must pass ruff lint (line-length 120, E/F/W rules).
- Process with Claude Code Opus 4.6 (1M token context).

TABLE REQUIREMENTS:
- Table 9: Global Regulatory Filing Matrix
  Columns: Jurisdiction | Filing Type | Timeline | Key Requirements | Agent Automation
  Include: FDA (IND), EMA (CTA/CTIS), PMDA (CTN), HC (CTA), TGA (CTN), NMPA (IND)
- Table 10: Accelerated Pathway Comparison
  Columns: Pathway | Jurisdiction | Eligibility | Benefit | Agent Decision Criteria

FORMATTING NOTES:
- Cross-reference \S\ref{sec:quality} for quality/compliance integration.
- Ensure all regulatory abbreviations are defined on first use.
- Reference the adapted 21 CFR 312 from the national platform paper.]

\subsection{Regulatory Agent}

[PROCESSING INSTRUCTION: Generate 1 page on the regulatory\_agent. This agent replaces
the Regulatory Affairs team. Cover: IND lifecycle management (initial IND filing per
21 CFR 312 Subpart B, serial amendments per 312.30, annual reports per 312.33),
authority correspondence (pre-IND meetings, Type A/B/C meetings), and multi-jurisdiction
coordination. Reference org\_01\_executive\_regulatory.md for sponsor accountability
under 21 CFR 312.50/312.52. Reference the adapted 21 CFR Part 312 from the
national platform. Process with Claude Code Opus 4.6 (1M token context).]

\subsection{IND and CTA Filing Automation}

[PROCESSING INSTRUCTION: Generate 1 page on filing automation. Cover: U.S. IND mechanics
(30-day review, safe-to-proceed), EU CTIS transition and ongoing management, UK CTIMP
reform, and eCTD v4.0 document assembly. Reference sponsor\_02\_preclinical\_regulatory.md
for IND/CTA mechanics. Reference sponsor\_06\_cmc\_submission.md for eCTD v4.0 and
submission readiness. Include CDISC SDTM/ADaM dataset standards.
Process with Claude Code Opus 4.6 (1M token context).]

\subsection{Accelerated Pathway Navigation}

[PROCESSING INSTRUCTION: Generate 1 page with Table 10 covering all accelerated pathways.
Reference sponsor\_02 for FDA formal meetings, Project Orbis, RTOR. Reference sponsor\_06
for global markets table (Health Canada NOC/c, Swissmedic fast-track, TGA provisional,
NMPA conditional). Describe how the autonomous system evaluates eligibility and
generates designation requests. Process with Claude Code Opus 4.6 (1M token context).]

\subsection{Regulatory Intelligence and Horizon Scanning}

[PROCESSING INSTRUCTION: Generate 0.5-0.75 page on automated regulatory monitoring.
Cover guidance update tracking, impact assessment algorithms, and compliance calendar
management. Reference the key regulatory drivers from sponsor\_01 (ICH E6(R3),
Project Optimus, Bayesian methods, EU CTR/CTIS, EU HTA).
Process with Claude Code Opus 4.6 (1M token context).]
